\section{Reference pipeline}
\label{sec:method}

The pipeline (Figure~\ref{fig:pipeline}) is deliberately conventional: a detector scores
edges, the highest-scoring edges seed candidate rings, and a ring layer grows and separates
them. It is the instrument for the measurements in Sections~\ref{sec:results}
and~\ref{sec:results-ceiling}, not the contribution.

\begin{figure}[t]
\centering
\begin{tikzpicture}[
  node distance=5mm and 6mm,
  stage/.style={draw, rounded corners=2pt, align=center, font=\small,
                minimum height=9mm, inner sep=5pt, fill=gray!8},
  arr/.style={-{Stealth[length=2mm]}, thick}
]
\node[stage] (l1) {Detector\\causal Multi-GINe\\edge scores};
\node[stage, right=of l1] (l2) {Seeds\\top 0.1\,\% of scores\\(chosen on validation)};
\node[stage, right=of l2] (l3a) {Expansion\\add edges of the\\same attempt};
\node[stage, right=of l3a] (l3b) {Partition\\one ring\\per attempt};
\node[stage, right=of l3b] (l3c) {Roles,\\pattern type};
\draw[arr] (l1) -- (l2);
\draw[arr] (l2) -- (l3a);
\draw[arr] (l3a) -- (l3b);
\draw[arr] (l3b) -- (l3c);
\end{tikzpicture}
\caption{The reference pipeline. ``B0'' denotes seeds grouped by connected components with no
further stage.}
% source: docs/report.md §3.2 (architecture, redrawn)
\label{fig:pipeline}
\end{figure}

\paragraph{Detector.}
A Multi-GINe edge classifier~\citep{egressy2024} with port numbering, time deltas, reverse
message passing, ego identifiers and subgraph-pattern edge features~\citep{blanusa2024},
trained with 5 seeds under batch-causal neighbour sampling (look-ahead within a batch of up to
1{,}800\,s; Section~\ref{sec:data}).
% source: docs/report.md §5.2
Causality matters here: a forward-looking variant of one cycle feature read edges after the
scored transaction and inflated its lift over the base rate by 12--16$\times$; made trailing,
the lift falls to 1.0--4.0$\times$ with primary-window edge F1 unchanged (0.495 vs.\ 0.494).
% source: docs/FINDINGS.md (C4)

\paragraph{Seeds.}
Seeds are the top $k^{*}=0.1\,\%$ of raw scores. The budget is fixed in advance for every detector
and seed; \texttt{min\_size} is chosen on validation ring-F1 (it is 3 throughout). As a sensitivity
analysis we also let validation choose $k^{*}$: it picks 0.1\,\% for every seed of Multi-GINe, but
0.05\,\% for two of five MEGA-GNN seeds, which lowers MEGA-GNN's full-pipeline ring-F1 from 0.602 to
0.510 (Table~\ref{tab:detector-swap}).
% source: configs/threshold_rule_k010.yaml; results/tables/r4_partition.csv (chosen_cells); results/tables/r6_detector_swap_sb2048_valpicked.csv (mega full 0.5098; seeds 4, 5 seed_frac_pct 0.05)

\paragraph{Expansion.}
A member edge below the seed cut --- a \emph{quiet member} --- never enters a baseline ring.
Expansion considers candidates that share an endpoint with an included edge within a
validation-chosen $\Delta t$ (up to two hops, capped at the 500 events nearest in time per
account to neutralise hub accounts).
% source: docs/EXPANSION.md §6.1
A LightGBM policy~\citep{ke2017} scores each candidate on 70 causal features: timing, money
conservation at the shared account, direction, reach to other seeds, trailing account state,
embedding similarity, detector-score quantile and trailing pattern counts.
% source: docs/EXPANSION.md §6.2
Its training data is cross-scored --- fold models trained on other days score each held-out
day --- so the policy never learns from in-sample detector scores; rows from day 0 are
dropped.
% source: docs/EXPANSION.md §6.3
The folds hold out days, not rings; Section~\ref{sec:limitations} reports the resulting leak
and its bound. We compare the policy with two rules: \emph{B1}, add every candidate sharing an
endpoint within $\Delta t$; and \emph{B1b}, the same with the candidate's score quantile above
0.99.
% source: docs/FINDINGS.md (B1)

\paragraph{Partition, roles and typing.}
A LightGBM affinity model scores pairs of ring edges that share an account, and
average-linkage clustering with a validation-chosen cut splits merged components. Each ring
then receives account roles from money-conservation rules and one of the eight pattern types
from a classifier over role counts, depth, cycle presence and size.
% source: docs/report.md §5.2

% Platform and explanation-layer paragraph moved to the released repository (cut rank 13).
